% =============================================================================
%  Tài liệu 06 - Báo cáo kết thúc dự án (Project Closure Report)
%  Build: latexmk 06-project-closure-report.tex   (pdfLaTeX, output in build/)
% =============================================================================
\documentclass[12pt,a4paper]{report}
\usepackage{sumtalk-report}
% =============================================================================
%  sumtalk-info.tex - information shared by every document (01-06).
%  Edit here once; every cover and footer picks it up.
% =============================================================================
\newcommand{\ProjectName}{SumTalk}
\newcommand{\CourseName}{QUẢN LÝ DỰ ÁN PHẦN MỀM}          % cover (upper case)
\newcommand{\CourseNameNormal}{Quản lý Dự án Phần mềm}     % footer
\newcommand{\AcademicYear}{2026 -- 2027}
\newcommand{\Instructor}{\canDien[họ tên giảng viên]}

% One \item per member: "Full name -- student ID", sorted by given name.
\newcommand{\TeamMembers}{%
  \item Nguyễn Anh Duy -- \canDien[MSSV]
  \item Đỗ Thành Đạt -- \canDien[MSSV]
  \item Trần Hải Đăng -- \canDien[MSSV]
  \item Châu Hùng Hào -- \canDien[MSSV]
  \item Bùi Đoàn Minh Nhật -- \canDien[MSSV]
  \item Võ Hoàng Sơn -- \canDien[MSSV]
}


\newcommand{\DocTitle}{BÁO CÁO KẾT THÚC DỰ ÁN}
\newcommand{\DocTitleNormal}{Báo cáo kết thúc dự án}
\newcommand{\DocTitleEn}{Project Closure Report}
\newcommand{\DocRunningTitle}{Báo cáo kết thúc dự án}
\newcommand{\DocCode}{SUMTALK-PCR}
\newcommand{\DocVersion}{2.1 -- 24/11/2026}
\newcommand{\DocStatus}{Chờ duyệt}
\newcommand{\DocAuthor}{\canDien[Project Manager]}
\newcommand{\DocReviewers}{Sponsor, Product Owner, Project Manager}
\newcommand{\DocReaders}{Sponsor, Product Owner, Project Manager}
\newcommand{\CoverDate}{TP. HỒ CHÍ MINH, THÁNG 11/2026}
\hypersetup{pdftitle={SumTalk -- Báo cáo kết thúc dự án},pdfauthor={Nhóm SumTalk}}

% "Passed" mark used in the assessment table.
\newcommand{\pass}{$\checkmark$}

\begin{document}

\makereportcover

% ============================================================ Front matter
\chapter*{Thông tin tài liệu}
\addcontentsline{toc}{chapter}{Thông tin tài liệu}

\docinfotable

\section*{Lịch sử phiên bản}

\begin{table}[H]
  \centering
  \begin{tabularx}{\linewidth}{@{}Q{1.6cm}Q{2.3cm}LP{3.2cm}@{}}
    \toprule
    \thead{Phiên bản} & \thead{Ngày} & \thead{Nội dung thay đổi} & \thead{Người thực hiện} \\
    \midrule
    1.0 & 12/11/2026 & Bản đầu, gồm ba tài liệu riêng: SOW và nghiệm thu, Release Notes, Lessons Learned
      & \canDien \\
    2.0 & 13/11/2026 & Gộp thành một tài liệu, đánh lại số mục (CR-08) & \canDien \\
    2.1 & 24/11/2026 & Chuyển sang định dạng báo cáo PDF (LaTeX), sắp xếp lại thành năm chương; thêm đánh giá sơ bộ
      theo mục tiêu M1--M6, dòng thời gian phát hành và biểu đồ phạm vi & \canDien \\
    \bottomrule
  \end{tabularx}
\end{table}

\needspace{8\baselineskip}
\heading{Quy ước trong tài liệu}
\begin{itemize}
  \item \canDien[\ldots]: số liệu đã có ở Azure Boards, dashboard chi phí hoặc biên bản họp nhưng chưa chép
        vào tài liệu, hoặc chỉ có khi kết thúc dự án; cần bổ sung trước khi trình duyệt.
  \item \textbf{Mốc M0--M5} và \textbf{mục tiêu M1--M6} theo điều lệ dự án~\cite{pid}. Hai hệ mã này dùng chung chữ M
        nhưng khác nhau: mốc gắn với thời điểm, mục tiêu gắn với kết quả.
  \item \textbf{Số liệu kỹ thuật} (độ trễ, thời gian xử lý, số test) lấy từ repo ngày 12/11/2026.
\end{itemize}

\tableofcontents

\chapter*{Tóm tắt}
\addcontentsline{toc}{chapter}{Tóm tắt}

Tài liệu này là \textbf{báo cáo kết thúc dự án} SumTalk: khép lại dự án bằng phạm vi đã thoả thuận và tiêu chí
nghiệm thu, các phiên bản đã phát hành, so sánh kế hoạch với thực tế, bài học kinh nghiệm và xác nhận đóng dự án.
Phần \emph{Statement of Work} (SOW) xác định tám hạng mục bàn giao, mỗi hạng mục có tiêu chí nghiệm thu kiểm tra
được, cùng quy trình xử lý thay đổi và mẫu biên bản nghiệm thu dùng cho từng mốc.

Bốn phiên bản đã được phát hành, từ v0.1 (ghi lại cuộc họp) tới \textbf{v0.4} ngày 12/11/2026 (app desktop,
transcript live dạng streaming, engine STT remote, bản web sẵn sàng deploy). So với kế hoạch, phạm vi tăng từ 4 lên
\textbf{7 epic} (thêm 29 story point); hiệu năng live tốt hơn mục tiêu (khoảng 1--2\,s so với ngưỡng 3--5\,s). Đánh
giá sơ bộ theo sáu mục tiêu cho thấy M2 và M4 đã đạt, M1, M3 và M5 đạt một phần, M6 chưa có số liệu. Dự án
\textbf{chưa đủ điều kiện đóng}: còn deploy thật lên AWS, chạy lại test E2E trình duyệt, UAT và nghiệm thu cuối.
Phần Lessons Learned ghi năm điều làm tốt, năm vấn đề cần cải thiện và năm bài học kỹ thuật đáng giữ.

\textbf{Từ khoá:} báo cáo kết thúc dự án, Statement of Work, nghiệm thu, release notes, lessons learned, kế hoạch
so với thực tế.

\chapter*{Danh mục từ viết tắt}
\addcontentsline{toc}{chapter}{Danh mục từ viết tắt}

\begin{center}
  \setstretch{1.1}
  \begin{tabular}{@{}P{2.2cm}P{13.2cm}@{}}
    \toprule
    \thead{Viết tắt} & \thead{Ý nghĩa} \\
    \midrule
    AC     & Actual Cost -- chi phí thực tế \\
    API    & Application Programming Interface -- giao diện lập trình ứng dụng \\
    AWS    & Amazon Web Services (EC2: máy chủ ảo; S3: lưu trữ đối tượng) \\
    BAC    & Budget at Completion -- ngân sách khi hoàn thành \\
    CI     & Continuous Integration -- tích hợp liên tục \\
    CPI    & Cost Performance Index -- chỉ số hiệu quả chi phí \\
    CR     & Change Request -- yêu cầu thay đổi \\
    E2E    & End-to-End -- kiểm thử đầu cuối \\
    GPU    & Graphics Processing Unit -- bộ xử lý đồ hoạ, dùng để chạy model \\
    HTTPS  & Hypertext Transfer Protocol Secure -- giao thức web có mã hoá \\
    IAM    & Identity and Access Management -- quản lý danh tính và quyền trên AWS \\
    LLM    & Large Language Model -- mô hình ngôn ngữ lớn \\
    NFR    & Non-Functional Requirement -- yêu cầu phi chức năng \\
    NSIS   & Nullsoft Scriptable Install System -- công cụ tạo bộ cài Windows \\
    PM     & Project Manager -- quản lý dự án \\
    PO     & Product Owner -- chủ sản phẩm \\
    POC    & Proof of Concept -- bản thử nghiệm khái niệm \\
    SLA    & Service Level Agreement -- cam kết mức dịch vụ \\
    SOW    & Statement of Work -- văn bản thoả thuận phạm vi công việc \\
    SP     & Story Point -- điểm ước lượng độ lớn của story \\
    SPI    & Schedule Performance Index -- chỉ số hiệu quả tiến độ \\
    STT    & Speech-to-Text -- chuyển giọng nói thành văn bản \\
    UAT    & User Acceptance Testing -- kiểm thử chấp nhận của người dùng \\
    VAC    & Variance at Completion -- chênh lệch ngân sách khi hoàn thành \\
    \bottomrule
  \end{tabular}
\end{center}

\cleardoublepage\phantomsection
\addcontentsline{toc}{chapter}{\listfigurename}
\listoffigures

\cleardoublepage\phantomsection
\addcontentsline{toc}{chapter}{\listtablename}
\listoftables

% ============================================================ Chapter 1
\chapter{Giới thiệu}
\label{chap:intro}

\section{Bối cảnh}

Dự án SumTalk được khởi động với sáu mục tiêu SMART và sáu mốc M0--M5~\cite{pid}, thực hiện theo Scrum qua bốn
sprint. Đến ngày 12/11/2026, Sprint 4 đã kết thúc với bản v0.4; các hồ sơ theo dõi (báo cáo sprint, Change Log)
nằm trong~\cite{mcr} và kết quả kiểm thử nằm trong~\cite{qtd}. Tài liệu này là bước cuối của vòng đời dự án: đối
chiếu những gì đã giao với những gì đã thoả thuận, và chuẩn bị cho việc nghiệm thu, đóng dự án.

\section{Mục đích và người đọc của tài liệu}

Tài liệu khép lại dự án ở năm khía cạnh: phạm vi đã thoả thuận và tiêu chí nghiệm thu, các phiên bản đã phát hành,
so sánh kế hoạch với thực tế, bài học kinh nghiệm, và xác nhận đóng dự án. Người đọc chính là Sponsor (bên giao việc,
nghiệm thu và xác nhận đóng), Product Owner (xác nhận sản phẩm đáp ứng backlog) và Project Manager (lập báo cáo,
chuyển bài học cho các dự án sau). Phần số liệu cuối dự án và bài học quản lý sẽ được điền khi kết thúc dự án.

\section{Cấu trúc tài liệu}

Phần còn lại của tài liệu gồm bốn chương:
\begin{itemize}
  \item \textbf{Chương~\ref{chap:sow} -- Statement of Work và nghiệm thu:} các bên, phạm vi và tiêu chí nghiệm thu,
        sản phẩm bàn giao, trách nhiệm của bên giao việc, quản lý thay đổi và mẫu biên bản nghiệm thu.
  \item \textbf{Chương~\ref{chap:release} -- Release Notes:} các phiên bản v0.1 đến v0.4 và mẫu cho phiên bản sau.
  \item \textbf{Chương~\ref{chap:lessons} -- Báo cáo tổng kết và Lessons Learned:} đánh giá theo mục tiêu, kế hoạch so
        với thực tế, điều làm tốt, điều cần cải thiện, bài học kỹ thuật và khuyến nghị.
  \item \textbf{Chương~\ref{chap:conclusion} -- Kết luận và xác nhận đóng dự án.}
\end{itemize}

% ============================================================ Chapter 2
\chapter{Statement of Work và nghiệm thu}
\label{chap:sow}

\textbf{SOW là thoả thuận giữa nhóm thực hiện và bên giao việc (Sponsor) về phạm vi giao, tiêu chí nghiệm thu và
cách xử lý thay đổi.} Tám hạng mục bàn giao đều có tiêu chí nghiệm thu dựa trên acceptance criteria của backlog,
yêu cầu phi chức năng hoặc test case đã có. Mẫu biên bản ở Mục~\ref{sec:acceptance-form} dùng cho từng mốc và cho
lần nghiệm thu cuối.

\section{Các bên}

Hai bên của SOW và vai trò của mỗi bên được nêu trong Bảng~\ref{tab:parties}.

\begin{table}[htbp]
  \centering
  \caption{Các bên của Statement of Work}
  \label{tab:parties}
  \begin{tabularx}{\linewidth}{@{}P{4.4cm}LP{4.2cm}@{}}
    \toprule
    \thead{Bên} & \thead{Đại diện} & \thead{Vai trò} \\
    \midrule
    Bên giao việc (khách hàng) & \canDien[đại diện Sponsor] & Phê duyệt phạm vi, nghiệm thu \\
    Bên thực hiện & Nhóm SumTalk -- \canDien[PM] & Xây dựng, bàn giao \\
    \bottomrule
  \end{tabularx}
\end{table}

\section{Mục đích}

Bên thực hiện xây dựng hệ thống SumTalk để ghi âm, phiên âm, tóm tắt và dịch cuộc họp, theo các mục tiêu trong
điều lệ dự án~\cite{pid}.

\section{Phạm vi công việc và tiêu chí nghiệm thu}

Bảng~\ref{tab:sow-scope} liệt kê tám hạng mục thuộc phạm vi. Tiêu chí nghiệm thu dẫn chiếu tới acceptance criteria
và yêu cầu phi chức năng trong đặc tả yêu cầu~\cite{srs}, và tới test case trong tài liệu chất lượng~\cite{qtd}.

\begin{table}[htbp]
  \centering
  \caption{Phạm vi công việc và tiêu chí nghiệm thu}
  \label{tab:sow-scope}
  \small
  \setlength{\tabcolsep}{4pt}
  \begin{tabularx}{\linewidth}{@{}Q{0.6cm}P{2.8cm}P{4.4cm}L@{}}
    \toprule
    \thead{\#} & \thead{Hạng mục} & \thead{Mô tả} & \thead{Tiêu chí nghiệm thu} \\
    \midrule
    1 & Ghi âm và upload & Trình duyệt, điện thoại (offline), upload $\leq$ 1\,GB
      & Acceptance criteria của US1.1--US1.3 đạt; test E2E trình duyệt đạt \\
    2 & Transcript & Ngôn ngữ gốc theo câu, live $\leq$ 3--5\,s, file $\geq$ 60 phút, xem, tìm, xuất
      & Acceptance criteria của US2.1--US2.4; NFR-P1, NFR-P2 \\
    3 & Tóm tắt & Story \#67 & 7 acceptance criteria của \#67; test case S-01\ldots S-11 \\
    4 & Dịch & Story \#60 & 6 acceptance criteria của \#60; test case T-01\ldots T-12 \\
    5 & App desktop & Bộ cài Windows & Acceptance criteria của US5.1; smoke test bản đóng gói 8/8 \\
    6 & Engine STT dùng chung & Remote, API \path{/v1} & Acceptance criteria của US6.1--US6.2 \\
    7 & Bản web trên AWS & Docker, HTTPS, đăng nhập, S3, hướng dẫn
      & Acceptance criteria của US7.1; checklist trong hướng dẫn deploy đạt trên AWS thật \\
    8 & Tài liệu & README, ARCHITECTURE, 6 tài liệu quản lý dự án, test case
      & Đủ file, không còn ô \canDien{} ở phần bắt buộc \\
    \bottomrule
  \end{tabularx}
\end{table}

\heading{Ngoài phạm vi}
Các hạng mục ngoài phạm vi giữ như trong điều lệ dự án~\cite{pid}: tài khoản và phân quyền, nhận diện người nói
(diarization), app iOS/Android native, và cam kết SLA thương mại.

\section{Sản phẩm bàn giao và thời hạn}

Năm sản phẩm bàn giao gắn với hai mốc cuối của dự án (Bảng~\ref{tab:sow-deliverables}): bộ cài và bản web thuộc mốc
M4, còn mã nguồn đã merge, tài liệu và báo cáo test thuộc mốc M5. Ngày của từng mốc theo điều lệ dự án~\cite{pid}.

\begin{table}[htbp]
  \centering
  \caption{Sản phẩm bàn giao và mốc}
  \label{tab:sow-deliverables}
  \begin{tabularx}{\linewidth}{@{}P{3.6cm}LQ{1.6cm}@{}}
    \toprule
    \thead{Sản phẩm} & \thead{Hình thức} & \thead{Mốc} \\
    \midrule
    Mã nguồn & Repo GitHub, các nhánh đã merge vào \path{main} & M5 \\
    Bộ cài Windows & \texttt{SumTalk-Setup-<version>.exe} & M4 \\
    Bản web & URL HTTPS và tài khoản test & M4 \\
    Tài liệu & Thư mục \path{docs/} & M5 \\
    Báo cáo test & Test Summary Report~\cite{qtd} & M5 \\
    \bottomrule
  \end{tabularx}
\end{table}

\section{Giả định và trách nhiệm của bên giao việc}

Để nhóm thực hiện bàn giao đúng hạn, bên giao việc cam kết ba điều:
\begin{itemize}
  \item \textbf{Phản hồi và nghiệm thu đúng hạn:} phản hồi trong Sprint Review; nghiệm thu trong vòng
        \canDien[3] ngày làm việc sau khi nhận bàn giao.
  \item \textbf{Chi phí hạ tầng:} cung cấp hoặc chấp nhận chi phí tài khoản AWS và API trong thời gian test, khoảng
        2,5--3 triệu VND/tháng với cấu hình CPU.
  \item \textbf{Dữ liệu đánh giá:} cung cấp mẫu audio thật (tiếng Việt, họp trộn ngôn ngữ) cho việc đánh giá chất
        lượng, nếu có.
\end{itemize}

\section{Quản lý thay đổi}

Mọi yêu cầu thay đổi đi theo quy trình của Change Log~\cite{mcr}. Thay đổi làm tăng hơn 10\% công sức hoặc dời mốc
phải được bên giao việc ký duyệt. Công sức phát sinh được tính theo đơn giá 100.000~VND/giờ, hoặc đổi lấy việc cắt
các hạng mục Should/Could có công sức tương đương.

\section{Mẫu biên bản nghiệm thu}
\label{sec:acceptance-form}

Hình~\ref{fig:acceptance-form} là mẫu biên bản dùng cho mỗi mốc M1--M5 và cho lần nghiệm thu cuối. Biên bản ghi rõ
phiên bản được nghiệm thu, kết quả theo từng hạng mục của Bảng~\ref{tab:sow-scope}, lỗi còn tồn và một trong ba kết
luận.

\begin{figure}[htbp]
  \centering
  \fbox{\begin{minipage}{0.94\linewidth}
    \setstretch{1.3}\small
    \begin{center}
      {\bfseries\large BIÊN BẢN NGHIỆM THU}\\[2pt]
      Mốc: M1\ldots M5 / nghiệm thu cuối \hfill Ngày: \dotfill
    \end{center}
    \textbf{Thành phần:} Bên giao việc: \dotfill \quad Bên thực hiện: \dotfill\par
    \textbf{Phiên bản nghiệm thu:} tag/commit, bộ cài, URL: \dotfill\par
    \medskip
    \renewcommand{\arraystretch}{1.5}
    \begin{tabularx}{\linewidth}{@{}P{3.2cm}LP{2.6cm}P{2.6cm}@{}}
      \toprule
      \thead{Hạng mục} & \thead{Tiêu chí} & \thead{Kết quả (Đạt / Không đạt)} & \thead{Ghi chú} \\
      \midrule
      & & & \\
      & & & \\
      & & & \\
      \bottomrule
    \end{tabularx}
    \medskip

    \textbf{Lỗi còn tồn} (mức, ID trên Azure Boards, hạn khắc phục): \dotfill\par
    \textbf{Kết luận:}\par
    \hspace*{1em}$\square$ Nghiệm thu toàn bộ\par
    \hspace*{1em}$\square$ Nghiệm thu có điều kiện (khắc phục trước ngày \makebox[2.6cm]{\dotfill})\par
    \hspace*{1em}$\square$ Không nghiệm thu\par
    \medskip
    \begin{tabularx}{\linewidth}{@{}CC@{}}
      \textbf{Bên giao việc} & \textbf{Bên thực hiện} \\
      (ký, ghi rõ họ tên) & (ký, ghi rõ họ tên) \\[1.6cm]
    \end{tabularx}
  \end{minipage}}
  \caption{Mẫu biên bản nghiệm thu}
  \label{fig:acceptance-form}
\end{figure}

% ============================================================ Chapter 3
\chapter{Release Notes}
\label{chap:release}

\textbf{Bốn phiên bản đã được phát hành, mỗi sprint một phiên bản; v0.4 ngày 12/11/2026 là bản mới nhất.} Mỗi
phiên bản liệt kê tính năng mới, thay đổi, lỗi đã sửa và hạn chế; chương trình bày bản mới nhất trước.
Hình~\ref{fig:releases} tóm tắt nội dung chính của từng phiên bản và bản v1.0 dự kiến sau nghiệm thu.

\begin{figure}[htbp]
  \centering
  \begin{tikzpicture}[
      ms/.style={circle,draw=white,line width=1pt,fill=primaryblue,inner sep=3pt},
      plan/.style={circle,draw=primaryblue,line width=1pt,fill=white,inner sep=3pt},
      lbl/.style={text width=2.9cm,align=center,font=\footnotesize,inner sep=1pt,
                  execute at begin node=\setstretch{1.05}}]
    % Solid: releases so far; dashed: planned v1.0.
    \draw[line width=1.2pt,primaryblue] (0,0) -- (12.0,0);
    \draw[-{Stealth[length=3mm]},line width=1.2pt,primaryblue,dashed] (12.0,0) -- (15.6,0);
    \foreach \x/\v/\when/\what in {
        1.4/v0.1/Sprint 1/{Ghi âm, offline, upload, transcript gốc, xuất file},
        7.4/v0.3/Sprint 3/{Tóm tắt (\#67), dịch (\#60), giao diện mới}} {
      \node[ms] at (\x,0) {};
      \node[lbl,above=7pt] at (\x,0) {\textbf{\v} -- \when\\ \what};
    }
    \foreach \x/\v/\when/\what in {
        4.4/v0.2/Sprint 2/{Transcript live, chọn model LLM},
        10.4/v0.4/{12/11/2026}/{App desktop, live streaming, STT remote, bản web sẵn sàng deploy}} {
      \node[ms] at (\x,0) {};
      \node[lbl,below=7pt] at (\x,0) {\textbf{\v} -- \when\\ \what};
    }
    \node[plan] at (13.6,0) {};
    \node[lbl,above=7pt] at (13.6,0) {\textbf{v1.0} -- dự kiến\\ Sau UAT và nghiệm thu cuối};
  \end{tikzpicture}
  \caption{Dòng thời gian các phiên bản phát hành}
  \label{fig:releases}
\end{figure}

\section{v0.4 -- desktop, transcript live dạng streaming, cloud}

\begin{sloppypar}
Phiên bản v0.4 thuộc Sprint 4, phát hành ngày 12/11/2026, gồm các nhánh \path{feature/desktop-app},
\path{feature/streaming-stt}, \path{feature/aws-deploy} và \path{docs/project-management}.
\end{sloppypar}

\heading{Mới}
\begin{sloppypar}
\begin{itemize}
  \item \textbf{App Windows (Electron):} mở như phần mềm, backend tự chạy và tự tắt; thu âm thanh hệ thống (Zoom,
        Teams bản desktop) mà không cần chia sẻ tab; bộ cài \texttt{SumTalk-Setup-0.3.0.exe} (khoảng 200\,MB).
  \item \textbf{Transcript live ``nói tới đâu hiện tới đó'':} từ mới hiện khoảng 1\,s sau khi nói; chữ đã chắc chắn
        in đậm, phần đang sửa in nghiêng mờ.
  \item \textbf{Engine STT remote} (\texttt{STT\_PROVIDER=remote}): dùng server GPU của nhóm, Groq hoặc OpenAI; máy
        khách không tải model.
  \item \textbf{API chuẩn OpenAI:} server có \path{/v1/audio/transcriptions}, bật bằng \texttt{STT\_SERVER\_API\_KEYS}.
  \item \textbf{File dài:} chia đoạn 5 phút tại chỗ im lặng; dòng đầu hiện sau vài giây; 1 giờ audio mất 1,5--5 phút
        trên GPU.
  \item \textbf{Bản web trên AWS:} Dockerfile, Docker Compose với Caddy (HTTPS tự động, đăng nhập), lưu audio và sao
        lưu database trên S3, hướng dẫn deploy từng bước.
  \item \textbf{Cấu hình sẵn trong bộ cài:} bộ cài có thể ghi sẵn cấu hình (\texttt{SUMTALK\_DEFAULT\_*}).
\end{itemize}
\end{sloppypar}

\heading{Thay đổi}
\begin{itemize}
  \item Chunk ghi âm giảm từ 3\,s xuống 1\,s; chu kỳ polling khi ghi từ 1,5\,s xuống 0,7\,s; cửa sổ live tối đa từ
        25\,s xuống 15\,s.
  \item Thêm thư viện \texttt{python-multipart} (server); \texttt{boto3} chỉ dùng cho bản server.
\end{itemize}

\heading{Hạn chế}
Bộ cài chưa ký số; bản GPU của app desktop là một thư mục khoảng 2,5\,GB (không đóng được bằng NSIS); bản web chưa
được deploy lên AWS thật.

\section{v0.3 -- tóm tắt và dịch}

Phiên bản v0.3 thuộc Sprint 3, trên nhánh \path{feature/sprint3-summary-translation}:
\begin{itemize}
  \item \textbf{Tóm tắt tự động} sau khi có transcript (\#67, \#68--\#71): tổng quan, ý chính, quyết định, việc cần
        làm, câu hỏi mở; tạo lại bằng tiếng Việt hoặc tiếng Anh; map-reduce cho transcript dài.
  \item \textbf{Dịch transcript} (\#60, \#61--\#65): chọn ngôn ngữ, thanh tiến độ, hiển thị song ngữ, xuất TXT/SRT.
  \item \textbf{Đánh giá và kiểm thử:} script đánh giá chất lượng (\#66, \#72); test E2E (\#99, \#100).
  \item \textbf{Giao diện mới:} sidebar, danh sách cuộc họp, recorder, player.
\end{itemize}

\section{v0.2 -- transcript live, chọn LLM}

Phiên bản v0.2 thuộc Sprint 2: \canDien[theo báo cáo Sprint 2: transcript live khoảng 3\,s, kết quả benchmark và
model LLM được chọn].

\section{v0.1 -- ghi lại cuộc họp}

Phiên bản v0.1 thuộc Sprint 1, gồm hai commit \texttt{c483e36 Initial commit} và \texttt{163e2fa Update the UI}:
ghi âm trên trình duyệt (micro, hoặc micro cùng âm thanh tab), ghi offline trên điện thoại, upload file lớn,
transcript ngôn ngữ gốc, xuất TXT/SRT/VTT/JSON/WAV.

\section{Mẫu cho phiên bản sau}

Các phiên bản sau dùng cùng một khung gồm sáu mục (Bảng~\ref{tab:release-template}), tiêu đề theo dạng
``vX.Y -- tên phiên bản (sprint, ngày phát hành)''.

% Last table of the chapter: [H] so it is not pushed onto a page of its own.
\begin{table}[H]
  \centering
  \caption{Mẫu release notes cho phiên bản sau}
  \label{tab:release-template}
  \begin{tabularx}{\linewidth}{@{}P{4cm}L@{}}
    \toprule
    \thead{Mục} & \thead{Nội dung} \\
    \midrule
    Nhánh / tag & Nhánh hoặc tag của phiên bản \\
    Mới & Tính năng mới \\
    Thay đổi & Thay đổi hành vi, cấu hình, thư viện \\
    Sửa lỗi & Lỗi đã sửa, kèm ID trên Azure Boards \\
    Hạn chế đã biết & Hạn chế còn lại của phiên bản \\
    Nâng cấp & Các bước cần làm khi cập nhật, ví dụ biến cấu hình mới \\
    \bottomrule
  \end{tabularx}
\end{table}

% ============================================================ Chapter 4
\chapter{Báo cáo tổng kết và Lessons Learned}
\label{chap:lessons}

\textbf{Đến Sprint 4, sản phẩm vượt mục tiêu về hiệu năng và mở rộng phạm vi từ 4 lên 7 epic, nhưng còn thiếu
deploy thật, UAT và số liệu chi phí để đóng dự án.} Phần kỹ thuật của chương được ghi từ kinh nghiệm đến hết
Sprint 4; phần số liệu và bài học quản lý sẽ được điền khi kết thúc dự án: \canDien.

\section{Đánh giá sơ bộ theo mục tiêu dự án}

Bảng~\ref{tab:objectives-check} đối chiếu sáu mục tiêu SMART của điều lệ dự án~\cite{pid} với bằng chứng đã có tại
ngày 12/11/2026 trong đặc tả yêu cầu~\cite{srs}, tài liệu kiểm thử~\cite{qtd} và báo cáo Sprint 4~\cite{mcr}. Đây
là đánh giá sơ bộ để chuẩn bị nghiệm thu, không thay cho biên bản nghiệm thu.

\begin{table}[htbp]
  \centering
  \caption{Đánh giá sơ bộ theo mục tiêu M1--M6 (12/11/2026)}
  \label{tab:objectives-check}
  \small
  \setlength{\tabcolsep}{4pt}
  \begin{tabularx}{\linewidth}{@{}Q{0.8cm}P{3.6cm}LP{2.3cm}@{}}
    \toprule
    \thead{Mã} & \thead{Mục tiêu} & \thead{Bằng chứng} & \thead{Đánh giá} \\
    \midrule
    M1 & Ghi âm hoặc tải file, transcript đúng ngôn ngữ người nói & Story E1, E2 Done; chạy trên Chrome/Edge, điện
      thoại, app Windows; test E2E trình duyệt cần chạy lại sau Sprint 4 & Đạt một phần (chờ E2E) \\
    M2 & Transcript hiện trực tiếp khi đang nói & Từ mới khoảng 1\,s, câu chốt trung bình 1,9\,s (mục tiêu
      $\leq$ 3--5\,s) & \pass{} Đạt \\
    M3 & Tóm tắt và dịch Việt $\leftrightarrow$ Anh & Story \#67, \#60 Done, test tự động đạt; 23 test case thủ công
      chưa có kết quả & Đạt một phần (chờ QC) \\
    M4 & Xử lý cuộc họp từ 60 phút trở lên & File 11,6 phút xong trong 18--57\,s, dòng đầu sau 3--5\,s; tóm tắt
      map-reduce cho transcript dài & \pass{} Đạt theo thước đo \\
    M5 & Bản web trên AWS và bộ cài Windows & Bộ cài v0.4 đạt smoke test 8/8; bản web chạy local, chưa deploy lên
      AWS thật & Đạt một phần \\
    M6 & Nằm trong ngân sách (CPI $\geq$ 0,9) & Chưa chép CPI từ dashboard & \canDien \\
    \bottomrule
  \end{tabularx}
\end{table}

Hai mục tiêu về hiệu năng (M2, M4) đã đạt với khoảng cách lớn. Ba mục tiêu M1, M3, M5 chỉ còn thiếu các bước kiểm
chứng cuối (E2E trình duyệt, test case QC, deploy thật), không còn thiếu tính năng; M6 cần số liệu chi phí.

\section{Kế hoạch so với thực tế}

Bảng~\ref{tab:plan-actual} so sánh kế hoạch với thực tế ở sáu khía cạnh. Phạm vi là khía cạnh thay đổi nhiều nhất.

\begin{table}[htbp]
  \centering
  \caption{Kế hoạch so với thực tế}
  \label{tab:plan-actual}
  \small
  \setlength{\tabcolsep}{4pt}
  \begin{tabularx}{\linewidth}{@{}P{2.4cm}P{2.9cm}P{2.9cm}P{1.9cm}L@{}}
    \toprule
    \thead{Khía cạnh} & \thead{Kế hoạch} & \thead{Thực tế} & \thead{Chênh lệch} & \thead{Nguyên nhân} \\
    \midrule
    Phạm vi & E1--E4 (ghi âm, transcript, tóm tắt, dịch) & E1--E7 (thêm desktop, engine dùng chung, bản web AWS)
      & +3 epic & Thay đổi CR-03\ldots CR-06 trong Change Log~\cite{mcr} \\
    Thời gian & \canDien[số] sprint & \canDien & & \\
    Chi phí nhân công (BAC) & \canDien & \canDien[AC] & VAC = & \\
    CPI / SPI cuối & 1,0 / 1,0 & \canDien & & \\
    Chất lượng & 0 lỗi Critical/High khi phát hành & \canDien & & \\
    Hiệu năng live & $\leq$ 3--5\,s & Khoảng 1--2\,s & Tốt hơn mục tiêu & Chunk 1\,s, LocalAgreement, engine GPU
      hoặc remote \\
    \bottomrule
  \end{tabularx}
\end{table}

Hình~\ref{fig:scope} lượng hoá phần tăng phạm vi theo story point đề xuất trong Product Backlog~\cite{srs}: bốn
epic ban đầu chiếm 105 SP; ba epic thêm vào (E5 app desktop, E6 engine STT dùng chung, E7 bản web) chiếm 29 SP,
khớp với tổng tác động của CR-03, CR-05 và CR-06 trong Change Log~\cite{mcr}. Như vậy phạm vi tăng khoảng 28\% so
với ban đầu.

\begin{figure}[htbp]
  \centering
  \begin{tikzpicture}
    \begin{axis}[
        width=0.86\linewidth, height=6.2cm,
        ybar, bar width=0.62cm, bar shift=0pt,
        ymin=0, ymax=40,
        symbolic x coords={E1,E2,E3,E4,E5,E6,E7},
        xtick=data,
        xticklabels={E1 Ghi âm,E2 Transcript,E3 Tóm tắt,E4 Dịch,E5 Desktop,E6 Engine STT,E7 Web},
        x tick label style={font=\footnotesize,text width=1.5cm,align=center},
        ylabel={Story point (đề xuất)},
        legend style={at={(0.97,0.95)},anchor=north east,font=\small,draw=gray!40},
        legend cell align=left,
        grid=major, xmajorgrids=false, grid style={gray!25},
        tick label style={font=\small},
        /pgf/number format/.cd, use comma, 1000 sep={.}]
      \addplot[fill=primaryblue, draw=primaryblue, nodes near coords,
               every node near coord/.append style={font=\footnotesize}]
        coordinates {(E1,34) (E2,34) (E3,18) (E4,19)};
      \addlegendentry{Phạm vi ban đầu (105 SP)}
      \addplot[fill=todored!70, draw=todored, nodes near coords,
               every node near coord/.append style={font=\footnotesize}]
        coordinates {(E5,8) (E6,13) (E7,8)};
      \addlegendentry{Thêm qua thay đổi (29 SP)}
    \end{axis}
  \end{tikzpicture}
  \caption{Phạm vi ban đầu và phần tăng thêm theo epic (story point đề xuất)}
  \label{fig:scope}
\end{figure}

\section{Điều làm tốt}

Năm cách làm sau đã giúp dự án giao được nhiều hơn kế hoạch mà vẫn giữ chất lượng, và nên được giữ cho các dự án
sau:
\begin{itemize}
  \item \textbf{Đo thay vì đoán:} độ trễ live và thời gian xử lý file dài được đo bằng script mô phỏng (số liệu trong
        Test Summary Report~\cite{qtd}), nên mọi tranh luận đều dựa trên số liệu.
  \item \textbf{Kiến trúc mở:} STT và LLM đều đi qua API chuẩn OpenAI. Đổi giữa Ollama, Groq, OpenAI hay server riêng
        chỉ bằng cấu hình; engine remote được thêm vào mà không sửa pipeline.
  \item \textbf{Không mất dữ liệu:} lưu chunk trên máy trước khi gửi; khôi phục sau crash; backend desktop tự tắt khi
        app chết.
  \item \textbf{Test tự động nhanh} (STT giả, LLM mock, S3 giả lập): 88 test chạy trong khoảng 6\,s, nên chạy được
        sau mỗi thay đổi.
  \item \textbf{Chia commit theo task trên nhánh feature:} dễ review, dễ đối chiếu với Azure Boards.
\end{itemize}

\section{Điều cần cải thiện}

Bảng~\ref{tab:improve} ghi các vấn đề đã gặp, tác động của chúng và bài học rút ra. Hai vấn đề đầu đều liên quan tới
việc quản lý mã nguồn, và cùng dẫn tới một rủi ro đã có trong Risk Register~\cite{pmp}: code chỉ nằm trên một máy
(R11).

\begin{table}[htbp]
  \centering
  \caption{Điều cần cải thiện và bài học}
  \label{tab:improve}
  \small
  \begin{tabularx}{\linewidth}{@{}P{4.4cm}P{3.6cm}L@{}}
    \toprule
    \thead{Vấn đề} & \thead{Tác động} & \thead{Bài học / hành động} \\
    \midrule
    Nhiều thay đổi Sprint 3 nằm chưa commit trong máy khi bắt đầu Sprint 4 & Khó tách commit, dễ mất code
      & Commit và push sau mỗi task; không để thay đổi qua đêm \\
    Push bị 403 do sai tài khoản hoặc quyền repo (R11) & Code chỉ có ở một máy, không review được
      & Kiểm tra quyền repo cho mọi thành viên từ Sprint 1; mỗi người dùng tài khoản của mình \\
    Phạm vi tăng 3 epic giữa dự án & Áp lực lịch và chi phí
      & Đánh giá tác động trước, cắt Should/Could tương ứng, ghi Change Log ngay \\
    Live trên CPU chậm (cửa sổ tới 25\,s được phiên âm lại mỗi lần) & Chữ hiện muộn trên máy không có GPU
      & Thiết kế live phải tính chi phí mỗi lượt; có engine remote ngay từ đầu cho máy yếu \\
    Bộ cài GPU quá lớn cho NSIS ($>$ 2\,GB) & Không có bộ cài GPU
      & Ưu tiên server GPU dùng chung thay vì GPU trên từng máy \\
    \canDien[bài học quản lý từ retro: họp, ước lượng, phân công] & & \\
    \bottomrule
  \end{tabularx}
\end{table}

\section{Bài học kỹ thuật đáng giữ}

Năm bài học kỹ thuật sau có giá trị vượt ra ngoài dự án này:
\begin{enumerate}
  \item Phiên âm live tốt cần ba thứ cùng lúc: chunk ngắn, cửa sổ ngắn (chỉ chốt các từ mà hai lượt phiên âm liên
        tiếp cho kết quả giống nhau), và đường truyền kết quả nhanh (polling ngắn hoặc push).
  \item Với họp trộn ngôn ngữ, nhận ngôn ngữ theo từng câu quan trọng hơn chọn model lớn hơn.
  \item Mọi API bên ngoài phải có retry, back-off và thông báo lỗi rõ trên giao diện; lỗi key phải hiện ngay ở header.
  \item Bản deploy POC chưa có tài khoản vẫn phải có HTTPS, đăng nhập, bucket private và IAM role tối thiểu.
  \item Môi trường dev có thể gây lỗi khó thấy. Ví dụ, VS Code đặt biến môi trường
        \texttt{ELECTRON\_RUN\_AS\_NODE=1}; vì vậy script khởi động phải tự làm sạch môi trường.
\end{enumerate}

\section{Khuyến nghị cho giai đoạn sau}

\begin{itemize}
  \item \textbf{Bảo mật:} làm tài khoản và phân quyền trước khi cho người ngoài dùng thật.
  \item \textbf{Tính năng:} nhận diện người nói; trang action items (\#79); dịch trực tiếp khi đang họp (\#84).
  \item \textbf{Phát hành:} ký số bộ cài, tự cập nhật app; CI chạy test tự động trên mỗi Pull Request.
  \item \textbf{Tài chính:} theo dõi chi phí API thật trong UAT để cập nhật Business Case~\cite{pid}.
\end{itemize}

% ============================================================ Chapter 5
\chapter{Kết luận và xác nhận đóng dự án}
\label{chap:conclusion}

\section{Tóm tắt}

Báo cáo đã chuẩn bị đầy đủ khung để nghiệm thu và đóng dự án SumTalk:
\begin{itemize}
  \item \textbf{SOW:} tám hạng mục bàn giao với tiêu chí nghiệm thu kiểm tra được, năm sản phẩm bàn giao gắn với mốc
        M4, M5, quy trình thay đổi và mẫu biên bản nghiệm thu.
  \item \textbf{Phát hành:} bốn phiên bản v0.1--v0.4; v0.4 mang lại app desktop, live streaming, engine STT remote và
        bản web sẵn sàng deploy.
  \item \textbf{Tổng kết:} M2, M4 đạt; M1, M3, M5 đạt một phần, chỉ còn các bước kiểm chứng; M6 chưa có số liệu.
        Phạm vi tăng 3 epic (khoảng 28\%), hiệu năng live tốt hơn mục tiêu.
  \item \textbf{Bài học:} năm điều làm tốt, năm vấn đề cần cải thiện, năm bài học kỹ thuật và bốn nhóm khuyến nghị.
\end{itemize}

\section{Xác nhận đóng dự án}

Dự án chỉ được đóng khi cả bốn hạng mục trong Bảng~\ref{tab:closure} được xác nhận.

\begin{table}[htbp]
  \centering
  \caption{Danh mục xác nhận đóng dự án}
  \label{tab:closure}
  \begin{tabularx}{\linewidth}{@{}LP{4.6cm}@{}}
    \toprule
    \thead{Hạng mục} & \thead{Trạng thái} \\
    \midrule
    Nghiệm thu cuối theo SOW (Chương~\ref{chap:sow}) & \canDien \\
    Bàn giao mã nguồn, bộ cài, URL, tài liệu & \canDien \\
    Gỡ tài nguyên AWS không dùng (EC2, Elastic IP, bucket) & \canDien \\
    Lưu trữ tài liệu, biên bản & \canDien \\
    \bottomrule
  \end{tabularx}
\end{table}

\section{Việc cần làm trước khi đóng dự án}

\begin{enumerate}
  \item Deploy bản web lên AWS thật và chạy checklist trong hướng dẫn deploy; chạy lại test E2E trình duyệt và bộ
        test Sprint 3; QC chạy 23 test case thủ công.
  \item Chạy UAT với người dùng thử, đo độ chính xác tiếng Việt; tổ chức nghiệm thu cuối bằng mẫu ở
        Hình~\ref{fig:acceptance-form}.
  \item Điền số liệu cuối dự án vào Bảng~\ref{tab:plan-actual} (số sprint, BAC, AC, CPI, SPI, lỗi khi phát hành) và
        mục M6 của Bảng~\ref{tab:objectives-check}.
  \item Bổ sung release notes v0.2 và bài học quản lý từ các buổi Retrospective.
  \item Điền tên đại diện hai bên, thời hạn nghiệm thu và các mục của Bảng~\ref{tab:closure}.
\end{enumerate}

\needspace{14\baselineskip}
\section{Phê duyệt}

Nhóm dự án đề nghị Sponsor, Product Owner và Project Manager xác nhận báo cáo kết thúc dự án này và ký đóng dự án
SumTalk sau khi các hạng mục của Bảng~\ref{tab:closure} hoàn tất.

\begin{table}[H]
  \centering
  \renewcommand{\arraystretch}{2.4}
  \begin{tabularx}{\linewidth}{@{}P{3.6cm}LP{3.2cm}P{2.6cm}@{}}
    \toprule
    \thead{Vai trò} & \thead{Họ tên} & \thead{Chữ ký} & \thead{Ngày} \\
    \midrule
    Sponsor & & & \\
    Product Owner & & & \\
    Project Manager & & & \\
    \bottomrule
  \end{tabularx}
\end{table}

% ============================================================ References
\begin{thebibliography}{9}
\addcontentsline{toc}{chapter}{\bibname}
  \bibitem{pid} Nhóm SumTalk, \emph{Hồ sơ khởi động dự án (Project Initiation Document)}, Tài liệu 01,
    mã SUMTALK-PID, phiên bản 2.1, 24/11/2026.
  \bibitem{srs} Nhóm SumTalk, \emph{Đặc tả yêu cầu phần mềm (Software Requirements Specification)},
    Tài liệu 02, mã SUMTALK-SRS, phiên bản 2.1, 24/11/2026.
  \bibitem{pmp} Nhóm SumTalk, \emph{Kế hoạch quản lý dự án (Project Management Plan)}, Tài liệu 03,
    mã SUMTALK-PMP, phiên bản 2.1, 24/11/2026.
  \bibitem{qtd} Nhóm SumTalk, \emph{Chất lượng và kiểm thử (Quality \& Test Document)}, Tài liệu 04,
    mã SUMTALK-QTD, phiên bản 2.1, 24/11/2026.
  \bibitem{mcr} Nhóm SumTalk, \emph{Báo cáo theo dõi và kiểm soát (Monitoring \& Control Report)}, Tài liệu 05,
    mã SUMTALK-MCR, phiên bản 2.1, 24/11/2026.
\end{thebibliography}

\end{document}
